
\subsection{同态的复合}
\label{subsec:4-6-3}

设 $L\subset M$ 是相对 spin 的拉格朗日子流形，
$\psi\colon M\to M'$ 是辛微分同胚。
在第 4.6.1 小节中，我们构造了从与 $L$ 关联的滤过
$A_\infty$ 代数到与 $L'=\psi(L)$ 关联的滤过
$A_\infty$ 代数的同态。
对另一个辛微分同胚 $\psi'\colon M'\to M''$，
也构造了从 $L'$ 的代数到
$L''=(\psi'\circ\psi)(L)$ 的代数的滤过
$A_\infty$ 同态。本小节证明，这两个同态的复合
与对应于 $\psi''=\psi'\circ\psi$ 的同态同伦
（定理 4.6.44）。

固定相容殆复结构 $J,J',J''$，
并取光滑奇异单形的三个可数集
\[
\mathcal X_a(L),\qquad
\mathcal X_b(L'),\qquad
\mathcal X_c(L'').
\]
它们分别张成 $C_a,C_b,C_c$。
在相应模空间上取多值截面
$\mathfrak s,\mathfrak s',\mathfrak s''$，
得到下列三个滤过 $A_\infty$ 代数：
\[
\begin{aligned}
&(C_a(L;\Lambda_{0,\mathrm{nov}}),\mathfrak m_k^{J,\mathfrak s}),\\
&(C_b(L';\Lambda_{0,\mathrm{nov}}),\mathfrak m_k^{J',\mathfrak s'}),\\
&(C_c(L'';\Lambda_{0,\mathrm{nov}}),\mathfrak m_k^{J'',\mathfrak s''}).
\end{aligned}
\]
由定理 4.6.1 与推论 4.6.2，有滤过 $A_\infty$ 同态
\[
\begin{aligned}
\mathfrak g^{\psi,\{J_\rho\}_\rho,
\mathfrak s^{\operatorname{top}(\rho)}}
&\colon C_a(L;\Lambda_{0,\mathrm{nov}})
 \longrightarrow C_b(L';\Lambda_{0,\mathrm{nov}}),\\
\mathfrak g^{\psi',\{J'_\rho\}_\rho,
\mathfrak s^{\operatorname{top}(\rho)'}}
&\colon C_b(L';\Lambda_{0,\mathrm{nov}})
 \longrightarrow C_c(L'';\Lambda_{0,\mathrm{nov}}),\\
\mathfrak g^{\psi'',\{J''_\rho\}_\rho,
\mathfrak s^{\operatorname{top}(\rho)''}}
&\colon C_a(L;\Lambda_{0,\mathrm{nov}})
 \longrightarrow C_c(L'';\Lambda_{0,\mathrm{nov}}).
\end{aligned}
\]
这里各源、靶均带上面所列的相应 $A_\infty$ 结构，
$\{J_\rho\}_\rho$ 从 $J^\psi$ 连到 $J'$，
$\{J'_\rho\}_\rho$ 从 $(J')^{\psi'}$ 连到 $J''$，
$\{J''_\rho\}_\rho$ 从 $J^{\psi''}$ 连到 $J''$。
下文把这三个同态及其构造中的
$\mathfrak f$ 分别简记为
$\mathfrak g^\psi,\mathfrak g^{\psi'},\mathfrak g^{\psi''}$
及 $\mathfrak f^\psi,\mathfrak f^{\psi'},\mathfrak f^{\psi''}$；
这些简记始终指向上面固定的路径与多值截面数据。
第 4.6.2 小节使用下标 $1,2,3$，
这里改用 $a,b,c$：构造 $\mathfrak g^\psi$ 时
$C_b(L';\Lambda_{0,\mathrm{nov}})$ 扮演 $C_2$，
构造 $\mathfrak g^{\psi'}$ 时它则扮演 $C_1$。

\begin{theorem}{4.6.44}\label{thm:4-6-44}
同态 $\mathfrak g^{\psi'}\circ\mathfrak g^\psi$
与 $\mathfrak g^{\psi''}$ 同伦。
\end{theorem}
\begin{proof}
上一小节已经证明 $\mathfrak g$ 的同伦类
不依赖于 $\psi$、$\{J_\rho\}_\rho$、
$\mathfrak s^{\operatorname{top}(\rho)}$ 等选择。
因此可取 $\{J''_\rho\}_\rho$ 为
\[
J''_\rho=
\begin{cases}
 (J_{2\rho})^{\psi'},&\rho\leq \frac12,\\
 J'_{2\rho-1},&\rho>\frac12.
\end{cases}
\]
这两支在 $\rho=\frac12$ 均为 $(J')^{\psi'}$。
\footnote{译者注：原书实体页 267 的两支分别印作
$(J'_{2\rho})^{\psi'}$ 与 $J''_{2\rho-1}$；
前者把已经在 $M''$ 上的路径再经 $\psi'$ 推前，
后者循环使用正在定义的路径。这里按本小节已定义的
两条路径及公共端点拼接。见疑误 A-0243。}
回顾 $\mathfrak g^{\psi''}$ 的构造使用下列模空间
\footnote{译者注：原书此处把环境流形印作 $M$；
但 $L''$ 与 $J''_\rho$ 均在 $M''$ 上，
下一个引理也以 $M''$ 为环境流形。
见疑误 A-0244。}：
\[
\mathcal M_{k+1}^{\mathrm{main}}
(M'',L'',\{J''_\rho\}_\rho:
\beta;\operatorname{top}(\rho)).
\]
关键是以下分解。

\begin{lemma}{4.6.45}\label{lem:4-6-45}
模空间
$\mathcal M_{k+1}^{\mathrm{main}}
(M'',L'',\{J''_\rho\}_\rho:
\beta;\operatorname{top}(\rho))$
同构于下列纤维积的不交并：
\[
\begin{aligned}
&\coprod_{\substack{\ell_1,\ldots,\ell_m\\
                    \sum_{i=1}^m\ell_i=k}}
 \coprod_{\beta_{\mathrm{last}}+\sum_i\beta_i=\beta}
\mathcal M_{m+1}^{\mathrm{main}}
(M'',L'',\{J'_\rho\}_\rho:
\beta_{\mathrm{last}};\operatorname{top}(\rho))\\[-1mm]
&\qquad {}_{L'}\!\times
\left(
\prod_{i=1}^m
\mathcal M_{\ell_i+1}^{\mathrm{main}}
(M',L',\{J_\rho\}_\rho:
\beta_i;\operatorname{top}(\rho))
\right).
\end{aligned}
\]
纤维积经 $\psi'\colon L'\to L''$ 识别评价点。
\end{lemma}
\begin{proof}
取
$((\Sigma,\vec z),(u_\alpha),(\rho_\alpha))$
为左端模空间的元素。
考虑满足 $\rho_\alpha>1/2$ 的分支，
并记这些分支的并为 $\Sigma_0$。
由 (4.6.7.2)，$\Sigma_0$ 连通。
将 $u,\rho$ 限制到 $\Sigma_0$，
得到
\[
\mathcal M_{m+1}^{\mathrm{main}}
(M'',L'',\{J'_\rho\}_\rho:
\beta_{\mathrm{last}};\operatorname{top}(\rho))
\]
的一个元素；其中 $m$ 是
$\Sigma\setminus\Sigma_0$ 的连通分支数，
这给出纤维积的第一个因子。
把其余连通分支记为
$\Sigma_i$（$i=1,\ldots,m$）。
将 $u,\rho$ 限制到各 $\Sigma_i$，
得到
$\mathcal M_{\ell_i+1}^{\mathrm{main}}
(M',L',\{J_\rho\}_\rho:
\beta_i;\operatorname{top}(\rho))$
中的元素。引理得证。
\end{proof}

图 4.6.6 展示这一分解中的一项。
\begin{figure}[!htbp]
 \centering
 \setlength{\unitlength}{0.67mm}
 \begin{picture}(118,105)
  \put(75,49){\circle{28}}
  \put(65,70){\circle{19}}
  \put(48,70){\circle{14}}
  \put(65,86){\circle{14}}
  \put(57,95){\circle{10}}
  \put(73,95){\circle{10}}
  \put(62,28){\circle{20}}
  \put(52,16){\circle{12}}
  \put(75,49){\makebox(0,0){$\scriptstyle\rho_0=0.8$}}
  \put(65,70){\makebox(0,0){$\scriptstyle0.6$}}
  \put(48,70){\makebox(0,0){$\scriptstyle0.2$}}
  \put(65,86){\makebox(0,0){$\scriptstyle0.3$}}
  \put(57,95){\makebox(0,0){$\scriptstyle0.1$}}
  \put(73,95){\makebox(0,0){$\scriptstyle0.1$}}
  \put(62,28){\makebox(0,0){$\scriptstyle0.4$}}
  \put(52,16){\makebox(0,0){$\scriptstyle0.3$}}
  \put(84,59){\line(1,1){3}}
  \put(72,77){\line(1,1){3}}
  \put(76,98){\line(1,1){3}}
  \put(53,99){\line(1,-1){3}}
  \put(54,36){\line(1,-1){3}}
  \put(69,20){\line(1,1){3}}
  \put(73,35){\line(1,0){4}}
  \put(85,62){\makebox(0,0)[l]{$\scriptstyle P_1$}}
  \put(75,80){\makebox(0,0)[l]{$\scriptstyle P_2$}}
  \put(80,103){\makebox(0,0)[l]{$\scriptstyle P_3$}}
  \put(51,103){\makebox(0,0)[r]{$\scriptstyle P_4$}}
  \put(52,39){\makebox(0,0)[r]{$\scriptstyle P_5$}}
  \put(73,18){\makebox(0,0)[l]{$\scriptstyle P_6$}}
  \put(76,31){\makebox(0,0)[l]{$\scriptstyle P_7$}}
  \put(89,49){\vector(1,0){11}}
  \put(102,49){\makebox(0,0)[l]{$\scriptstyle ev_0$}}
 \end{picture}
 \par\smallskip
 \renewcommand{\thefigure}{4.6.6}
 \refstepcounter{figure}
 \textbf{图~\thefigure}\label{fig:4-6-6}\quad
 $\mathfrak f^{\psi'}_6
 \bigl(\psi(P_1),\psi(P_2),
 \mathfrak f^\psi_2(P_3,P_4),
 \mathfrak f^\psi_0(1),
 \mathfrak f^\psi_2(P_5,P_6),\psi(P_7)\bigr)$
\end{figure}

可选取与引理 4.6.45 的同构相容的
多值截面 $\mathfrak s^{\operatorname{top}(\rho)''}$，
于是
\begin{equation}
\mathfrak f^{\psi'}\circ\mathfrak f^\psi
=\mathfrak f^{\psi''}.
\tag{4.6.46}\label{eq:4-6-46}
\end{equation}
更准确地说，有图式 4.6.3 所示的交换图。
若采用第 4.6.2 小节的记号，其中的
$C_{a,b}(L';\Lambda_{0,\mathrm{nov}})$
就是 $C_3(L';\Lambda_{0,\mathrm{nov}})$。
图中为清晰起见省略各复形共有的系数
$\Lambda_{0,\mathrm{nov}}$；
带下标 $a,c$、$b,c$、$(a,b),c$ 的复形
均以 $L''$ 为底空间。
\footnote{译者注：原书实体页 268 把图中
$C_{(a,b),c}$ 及下文扩大复形的底空间写作
未定义的 $L'_1$；按两侧子复形与同态靶空间，
应为 $L''$。见疑误 A-0246。}
\begin{figure}[htbp]
 \centering
 \setlength{\unitlength}{0.72mm}
 \begin{picture}(162,89)
  \put(15,75){\makebox(0,0){$C_{a,c}(L'')$}}
  \put(143,75){\makebox(0,0){$C_c(L'')$}}
  \put(79,46){\makebox(0,0){$C_{(a,b),c}(L'')$}}
  \put(143,46){\makebox(0,0){$C_{b,c}(L'')$}}
  \put(15,12){\makebox(0,0){$C_a(L)$}}
  \put(79,12){\makebox(0,0){$C_{a,b}(L')$}}
  \put(143,12){\makebox(0,0){$C_b(L')$}}
  \put(130,75){\vector(-1,0){99}}
  \put(143,68){\vector(0,-1){15}}
  \put(128,46){\vector(-1,0){26}}
  \put(15,19){\vector(0,1){48}}
  \put(26,12){\vector(1,0){37}}
  \put(129,12){\vector(-1,0){35}}
  \put(79,19){\vector(0,1){19}}
  \put(143,19){\vector(0,1){19}}
  \put(30,68){\vector(2,-1){36}}
  \put(18,44){\makebox(0,0)[l]{$\scriptstyle\mathfrak f^{\psi''}$}}
  \put(42,15){\makebox(0,0)[b]{$\scriptstyle\mathfrak f^\psi$}}
  \put(81,28){\makebox(0,0)[l]{$\scriptstyle\mathfrak f_{\mathrm{ext}}^{\psi'}$}}
  \put(145,28){\makebox(0,0)[l]{$\scriptstyle\mathfrak f^{\psi'}$}}
 \end{picture}
 \par\smallskip
 \renewcommand{\thefigure}{4.6.3}
 \refstepcounter{figure}
 \textbf{图式~\thefigure}\label{diag:4-6-3}
\end{figure}

依定义，由 $C_a(L)$ 经 $\mathfrak f^\psi$
到 $C_{a,b}(L')$，再经自然包含
$C_b(L')\hookrightarrow C_{a,b}(L')$
的同伦逆到 $C_b(L')$ 的复合，
就是 $\mathfrak g^\psi$。
同样，图式第三列的两个箭头复合，
即把其中一个箭头的方向用同伦逆反转，
得到 $\mathfrak g^{\psi'}$。
第一列的箭头再与第一行自然包含的同伦逆复合，
得到 $\mathfrak g^{\psi''}$。

由 (4.6.46)，可取扩大复形
\[
C_{(a,b),c}(L'';\Lambda_{0,\mathrm{nov}})
\supset
C_{a,c}(L'';\Lambda_{0,\mathrm{nov}})
\cup C_{b,c}(L'';\Lambda_{0,\mathrm{nov}})
\]
和延拓同态 $\mathfrak f_{\mathrm{ext}}^{\psi'}$，
使图式 4.6.3 交换（比较引理 7.2.299）。
于是定理 4.6.44 得证。
\end{proof}

还可以把定理 4.6.44 中的“同伦”
加强为“hu-同伦”。这一点可由上面的论证
与第 7.3 节的论证合并得到。
